\xchapter{基于物理模拟与深度学习的雷达信号仿真及迁移学习优化}{?}
\xsection{引言}{Introduction}
在上一章中，本文提出了一种基于超宽带虚拟阵列的物体三维点云重建方法，该方法通过为超宽带雷达构建出大量的虚拟天线来提升雷达感知能力，并通过背景反射消减和有效距离箱定位提取出有用的雷达信号，之后通过设计的神经网络模型从雷达信号中恢复出物体多个视角下的深度图，并最终通过对多视角深度图进行逆投影和点云融合得到了物体最终的重建点云。然而由于该方法使用深度学习的方法从雷达信号去恢复深度图，要想训练出足够性能的网络模型依赖于大量的带多视角深度图标注的雷达信号数据，然而真实数据的采集受限于信号采集复杂程度、多视角深度图标注难度等，这导致了训练数据稀缺且多样性不足；同时，神经网络模型性能与训练数据所处数据域高度绑定，当将网络模型迁移到其他场景下的新的数据域上时，模型性能会显著下降，此时需要在新数据域上从头开始训练，这会消耗大量的时间以及算力成本。

针对上述问题，本章介绍本文的第二个贡献点，即提出了一种基于物理模拟与深度学习的雷达信号仿真方法，该方法能针对公开数据集中大量已有的物体三维模型生成对应的仿真雷达信号数据，同时这些公开数据集中的物体三维模型容易去渲染生成多视角深度图，将仿真雷达信号数据与对应的多视角深度图标注组成一个数据量庞大的仿真数据集。之后采用迁移学习中的“预训练(Pre-training)-微调(Fine-tuning)”两阶段学习范式：首先将网络模型在仿真数据集上进行充分的预训练，得到从仿真数据集中学习到通用知识的预训练模型(Pre-trained Model)，之后使用预训练模型的模型参数在真实数据集上进行简单微调即可完成模型的迁移。在预训练模型上进行微调相比于从头训练的优势在于：（1）数据效率高，预训练模型已经从大规模数据中学习到了通用知识（雷达信号的特征，以及从特征到物体三维结构的映射关系），因此，只需要使用少量的真实数据在预训练模型上微调就能取得良好的性能表现。（2）训练成本低，相比于随机初始化的参数，预训练模型的参数已经接近于最优解，这大大减少了模型需要搜索的参数空间范围，因此只需要少量迭代即可让模型在真实数据集上收敛。（3）模型性能更好，由于对预训练模型高质量特征表示的直接继承，以及更稳定的训练优化过程，微调后的模型相比于从头训练的模型具有更好的性能表现。（4）更便捷的跨域迁移能力，得益于使用预训练模型在真实数据集上微调时极高的数据效率和极低的训练成本，使得模型可以快速迁移到其他场景下的数据域上。

\xsection{基于物理模拟的雷达信号生成}{Physics-based Radar Signal Generation}
本小节探究如何通过物理模拟的方法根据物体三维模型生成对应的初步仿真雷达信号。一些三维重建领域公开的数据集中包含了大量的物体三维模型，例如ShapeNet数据集\cite{chang2015shapenet}，它是一个大规模、多类别的三维模型数据库，其中给出了每个三维模型高精度的多种三维表示，例如点云、三维网格等，本文通过物理建模的方式模拟超宽带信号经过三维模型反射后到达雷达接收端的接收信号，进而为大量三维模型生成对应的仿真雷达信号。
\xsubsection{三维物体反射信号物理建模}{?}
在章节\ref{sec:UWB_sense}中已经对雷达接收信号进行了详细的数学建模，具体见公式\ref{equ:y_complex}，然而在该建模中，最终的接收信号被简单表示成多个路径反射信号的叠加，但并没有对每条路径的物理意义进行进一步讨论，如果要对三维物体的反射信号进行精确建模，必须考虑物体可能产生的每条反射路径。根据文献\cite{korany2019xmodal}中对人体反射WiFi信号的建模方式，人体的整体反射信号可以被看做是人体三维网格(Mesh)的每个三角形面(Face)的反射信号的叠加，这种建模方式细化了人体作为反射物时产生的反射路径的粒度，本文参考该文献中的建模方式，对三维物体的超宽带反射信号进行了物理模拟，如图{?}所示。首先将反射路径细化为三维网格的每个三角形面，对雷达接收信号进行重写：
\begin{equation}
        y(t) = \sum_{f=1}^{F}A_{f}\cdot e^{j\theta_f} + n(t)
\end{equation}
其中三角形面$f$对应的反射信号幅值为$A_{f} = \frac{1}{2}\alpha_fg(t-kT_s-T_f)$，相位为$\theta_f$为$2\pi f_c T_f$。

在之前的讨论中，认为$\alpha_p$是反射路径$p$的固有反射强度系数，这里，$\alpha_f$也可以认为是三角形面$f$的固有反射强度系数，为了准确估计每个三角形面的反射信号幅值，需要对$\alpha_f$进行进一步的详细数学建模。对于每个三角形面的反射信号强度，会受到三个因素的影响：即